\chapter{A Primer on Latency, Tokens, and Context}
\label{ch:primer}

\epigraph{``I feel the need. The need for speed.''}{Maverick and Goose,
\emph{Top Gun} (1986)}

Everybody quotes that line. Almost nobody remembers what happens next, which is
that the two of them proceed to fly recklessly, get chewed out, and eventually
lose an aircraft. The need for speed is not a strategy. It is an appetite. What
turns it into engineering is knowing which knob actually moves the aircraft.

So before we look at a single byte of the Model Context Protocol, we are going
to spend a chapter on the physics. Not because physics is fun, though it is.
Because every design decision in the rest of this book is a trade against one of
exactly three budgets. Miss that, and you will spend three weeks optimising the
wrong thing. I have done this. I have
watched a very good team spend a sprint moving a server from Python to Go, ship
it, measure it, and discover they had improved end-to-end task latency by
approximately one half of one percent.

They were not stupid. They were optimising the only layer they could see.

\section{The three budgets}
\label{sec:three-budgets}

Every agentic application spends three things. Not two. Not four. These three,
and they are coupled, so you cannot economise on one without paying in another.

\bookfig[0.80]{three-budgets}{The three budgets, and the fact that they trade
against each other. Almost every optimisation in this book is a move along one
of these edges. The trick is knowing which direction you are moving.}{three-budgets}

\subsection{Latency}

Wall clock, from the user pressing enter to the user having an answer. It
decomposes into pieces you must learn to name separately, because they behave
completely differently:

\begin{description}[leftmargin=1.4em, style=nextline, itemsep=4pt]
\item[Time to first token (TTFT)] How long the model thinks before it emits
anything. Dominated by prompt processing, so it grows with how much context you
sent. This is the number your users feel as ``is it broken?''
\item[Tokens per second] The generation rate once it starts. Roughly constant
for a given model, and roughly outside your control.
\item[Round-trip time of one loop iteration] Model turn, plus tool call, plus
the model reading the result. This is the unit that actually matters, and it is
the one nobody measures.
\end{description}

The third one is the whole ballgame. A single agent task is not one model call.
It is a loop, and every trip round that loop pays for a full model turn.

\subsection{Context}

The token window is your bandwidth. It is also your memory, which is an unusual
and slightly horrible property: the same finite resource stores the
conversation, the tool catalogue, the system prompt, and every result you have
retrieved so far. Fill it with junk and there is no room for the task.

Worse, it is not a one-time cost. Anything sitting in context is re-sent and
re-processed on every single turn. A tool description you added in March is
being paid for, right now, on a request that will never call it.

\bookfig[0.86]{xkcd-context}{The context window of a working agent, in
approximate proportion. The green sliver is the part anybody would defend in a
code review.}{xkcd-context}

\subsection{Cost}

Dollars per task, and it deserves to be a first-class metric rather than
something finance discovers in the quarterly. It is nearly a linear function of
context: input tokens times input price, plus output tokens times output price,
summed over every turn in the loop.

Which means context bloat is not merely a memory problem. It is a recurring
bill, and it scales with your traffic, which is exactly the direction you were
hoping your traffic would scale.

\keyidea{Latency, context, and cost are one resource wearing three hats. Every
optimisation in this book is a trade among them, and most of the good ones
attack all three at once because they remove work rather than moving it.}

\section{Anatomy of one loop iteration}
\label{sec:anatomy}

Here is where the time goes. This is the diagram the rest of the book keeps
returning to, so it is worth staring at for a minute.

\bookfig[0.94]{loop-anatomy}{One iteration of an agent loop, drawn to scale.
The model bands are simulated from a stated distribution; the server band is
measured. Note the proportion, and note which band this book can actually help
you with.}{loop-anatomy}

Read it left to right:

\begin{enumerate}[leftmargin=1.6em, itemsep=3pt]
\item \textbf{The host assembles context.} Conversation history, system prompt,
and the merged tool catalogue from every connected server. Cheap in time.
Ruinously expensive in tokens, which we will get to.
\item \textbf{The model thinks.} Time to first token, then generation. It emits
a tool-call intent: a name and some arguments.
\item \textbf{The host serialises and sends.} JSON-RPC, over stdio or HTTP.
Microseconds of CPU, plus whatever the network costs.
\item \textbf{The server executes.} Your actual business logic. The part you
own, the part you profile, and usually the smallest slice on the chart.
\item \textbf{The result comes back and re-enters context.} And now the model
has to read it, which means the whole prompt gets processed again.
\item \textbf{The model thinks again.} Another full turn. This is step 2 again,
and it costs what step 2 cost.
\end{enumerate}

Steps 5 and 6 are the expensive ones and they are the ones people forget,
because they do not look like work. Nothing is executing. But a tool result
entering context means the next model turn processes a longer prompt, and you
pay another TTFT, and you are billed for every token again.

\subsection{What the proportions actually are}

Enough theory. Meridian, the reference application that ships with this book,
runs a three-step task against four MCP servers: score an account for credit
risk, screen it for fraud, and price a facility against the reference curve.
Here is where its wall clock goes.

\begin{measurebox}{Meridian baseline, one task, three iterations}
\begin{tabular}{@{}lrr@{}}
\toprule
\thead{Band} & \thead{Milliseconds} & \thead{Share} \\
\midrule
Model inference (3 turns)          & 1{,}312.5 & 96.1\,\% \\
Transport + server execution       & 53.9      & 3.9\,\% \\
\midrule
\textbf{Total wall clock}          & \textbf{1{,}366.4} & \\
\bottomrule
\end{tabular}

\vspace{6pt}
\footnotesize
3 model turns, 3 tool calls, 4 servers, 4{,}009 tokens, \$0.0125 per task.
Model timings simulated from the distribution in
\secref{simulated}; transport and execution measured on the machine in
\secref{machine}. Regenerate with \texttt{make bench}.
\end{measurebox}

\bookfig[0.92]{loop-breakdown}{The same numbers, to scale. If you have been
optimising the green sliver, this chart is either very good news or very bad
news depending on how long you have been at it.}{loop-breakdown}

Now, before you conclude that the protocol does not matter and close the book:
that green sliver is not the point. The point is the \emph{number of blue
bands}. There are three of them because the loop ran three times.

\keyidea{You do not make an agent fast by making the tool call fast. You make
it fast by needing fewer tool calls, and by making each one carry more.}

That reframing is the reason this book exists. Every technique in
Chapter~\ref{ch:optimizing} is ultimately a way of deleting a blue band, or of
stopping one from getting wider.

\section{Round trips are the enemy}
\label{sec:round-trips}

If you take one idea from this chapter, take this one.

An extra round trip is not an extra 20 milliseconds of network time. It is an
extra \emph{model turn}, because the model has to look at what came back and
decide what to do next. On the Meridian baseline that is roughly 440
milliseconds, plus another 1,300 tokens of input, plus the cost of those tokens,
every time.

Watch what that does to an MRTR interaction. Multi round-trip requests are how a
2026-era server asks the client a question mid-call, and Chapter~\ref{ch:mrtr}
covers the mechanism properly. For now, only the price tag matters.

\begin{measurebox}{The cost of one elicitation}
\begin{tabular}{@{}lrr@{}}
\toprule
\thead{} & \thead{Requests} & \thead{Latency} \\
\midrule
Tool call, no elicitation           & 1.0 & 34.0\,ms \\
Tool call with elicitation          & 2.0 & 68.5\,ms \\
\midrule
Transport delta                     &     & +34.5\,ms \\
Plus one model turn to read the form&     & +340\,ms \\
\textbf{True delta}                 &     & \textbf{+375\,ms} \\
\bottomrule
\end{tabular}

\vspace{6pt}
\footnotesize
Measured at 12\,ms simulated per-call latency. The transport cost of the extra
round trip is real but small; the model turn it forces is ten times larger.
\end{measurebox}

The transport delta is 34 milliseconds and looks harmless. The honest number is
375, because a round trip drags a model turn along behind it. This is why
Chapter~\ref{ch:mrtr} spends so much energy on collecting inputs up front, and
why ``just add a confirmation prompt'' is a more expensive suggestion than it
sounds in a design review.

The same arithmetic runs the other way, which is the good news. When the model
asks for three tools at once, it is telling you those three are independent.
Run them together:

\begin{measurebox}{Serial versus parallel fan-out, three independent tools}
\begin{tabular}{@{}rrrr@{}}
\toprule
\thead{Per-call latency} & \thead{Serial} & \thead{Parallel} & \thead{Speedup} \\
\midrule
2\,ms   & 18.6\,ms  & 7.0\,ms  & 2.7$\times$ \\
10\,ms  & 85.8\,ms  & 30.5\,ms & 2.8$\times$ \\
40\,ms  & 286.7\,ms & 98.1\,ms & 2.9$\times$ \\
\bottomrule
\end{tabular}

\vspace{6pt}
\footnotesize
Three calls, so the ceiling is 3$\times$. The gap between 2.8 and 3.0 is thread
scheduling and the slowest of the three servers.
\end{measurebox}

Notice that the speedup gets \emph{better} as latency rises. Serial fan-out
pays the sum of the latencies; parallel pays the maximum. The worse your
network, the more the mistake costs you, which is a deeply unfair property of
the universe and one you should exploit.

\section{The N times M problem, and why a protocol is the answer}

Let us back up and ask why MCP exists, because the answer is not
``standardisation is nice.''

You have $N$ AI applications and $M$ systems worth connecting them to. Without a
shared contract, you write $N \times M$ integrations, and worse, you
\emph{maintain} $N \times M$ integrations. Four hosts and five systems is twenty
pieces of bespoke glue, each with its own auth, its own error handling, and its
own person who understands it and is currently on holiday.

\bookfig[0.98]{nxm-problem}{Twenty integrations, or nine. This is the same
argument that produced ODBC, LSP, and every other narrow waist in computing
history, and it works for the same reason.}{nxm-problem}

With a protocol you write $N + M$. Nine. Each host implements a client once,
each system exposes a server once, and the contract in the middle is somebody
else's problem to specify.

This is the Language Server Protocol argument, and it is worth dwelling on why
LSP worked, because it tells you what MCP has to get right. Before LSP, every
editor implemented Go To Definition for every language. After LSP, editors
implemented one client and languages implemented one server, and the ecosystem
exploded. The waist was narrow enough to implement in a weekend and expressive
enough to be worth implementing.

\subsection{Why not a framework?}

Fair question, since frameworks got there first.

A framework is a library you import. It works beautifully as long as everyone
is in the same process, in the same language, on the same release cadence. The
moment your risk team ships Python, your platform team ships TypeScript, and a
vendor ships a binary you cannot recompile, a library is the wrong shape. You
need something that survives a process boundary and a version skew.

That is a protocol. And once it is a protocol, it needs a wire format, a
versioning story, and a deprecation policy, which is precisely what
Chapter~\ref{ch:architecture} is about.

\subsection{Positioning against the alternatives}

You have choices, and MCP is not always the right one. Here is my honest read as
of mid-2026.

\begin{table}[htbp]
\centering
\small
\begin{tabularx}{\textwidth}{@{}lXX@{}}
\toprule
\thead{Approach} & \thead{Optimises for} & \thead{Costs you} \\
\midrule
Provider-native function calling
  & Lowest possible friction; zero extra infrastructure
  & Rewrite per provider; no discovery, consent, or resource story \\
OpenAPI-driven tool use
  & Reusing an API surface you already document
  & REST shape is wrong for planning; catalogues balloon (\chapref{tools}) \\
Framework tool abstractions
  & Speed inside one language and one process
  & Version lock-in; no cross-organisation reuse \\
\textbf{MCP}
  & \textbf{Many hosts, many systems, separate teams, real trust boundaries}
  & \textbf{A protocol to learn, and a host to run} \\
A2A and agent-to-agent protocols
  & Agents negotiating with peer agents
  & Different problem; composes with MCP rather than replacing it \\
\bottomrule
\end{tabularx}
\caption{Where each approach earns its keep. If you are one team, one language,
one model provider, and one system, provider-native function calling is
genuinely the right answer and you should use it.}
\label{tab:positioning}
\end{table}

The A2A comparison is the one people ask about most, so: they are orthogonal.
MCP connects an agent to \emph{capabilities}. A2A connects an agent to
\emph{other agents}. Chapter~\ref{ch:multi-agent} builds a system that is both,
because an agent that consumes MCP tools on one side and exposes MCP tools on
the other is a completely ordinary thing to want.

\section{Meet Meridian}
\label{sec:meridian}

Meridian is the instrumented reference application for this book. It is a
financial analytics assistant, and it exists so that no claim in these pages has
to be taken on faith.

\begin{description}[leftmargin=1.4em, style=nextline, itemsep=3pt]
\item[\texttt{meridian-risk}] Credit scoring. Tools, resources, prompts, MRTR
approvals, the Tasks extension, and an MCP App template. The busiest server.
\item[\texttt{meridian-compliance}] Transaction review with a human in the loop,
an immutable audit log, and (on request) the tool-poisoning payload from
Chapter~\ref{ch:threats}.
\item[\texttt{meridian-fraud}] Fast, read-only behavioural signals. The cheap
parallel call.
\item[\texttt{meridian-marketdata}] Reference rates. Everything it returns is
identical for every caller, which makes it the one place a shared gateway cache
genuinely pays.
\end{description}

Plus a host, a client, an agent loop, a measurement harness, and 134 tests.

The whole thing has no third-party runtime dependencies, which was a deliberate
and slightly obsessive decision. The SDKs are good and you should use them in
production. But they were still shipping the handshake-era protocol while this
book was being written, and a book about the wire that hides the wire behind
somebody's convenience wrapper is not doing its job. So \texttt{meridian/protocol/}
is a complete implementation of \texttt{2026-07-28} in about two thousand lines
of standard-library Python. You can read all of it in an evening.

\subsection{Try it now}

\begin{sh}
git clone https://github.com/krimler/Building-Awesome-MCP-Apps
cd Building-Awesome-MCP-Apps
make test     # 134 tests, roughly two seconds
make bench    # every number in this book
\end{sh}

\noindent It also connects to real clients. This is Claude Code driving three
Meridian servers at once:

\begin{sh}
claude -p "Assess ACC-1042 for risk, screen it for fraud, and get the 1Y and
5Y reference rates." --mcp-config .mcp.json
\end{sh}

\begin{plain}
RISK:  55.4 elevated
FRAUD: watch 1
CURVE: 1Y=3.96 5Y=3.68
\end{plain}

\noindent Those values match the fixtures exactly, which is the point.
\texttt{meridian/VERIFICATION.md} records the full transcript.

There is a wrinkle in that command worth flagging now, because it shaped several
later chapters. Claude Code, like most clients shipping in 2026, still opens a
connection with an \texttt{initialize} handshake. Revision \texttt{2026-07-28}
deleted that handshake. A strictly conforming modern server would therefore be
completely correct and completely unusable. Meridian ships a dual-era bridge to
handle exactly this, and Chapter~\ref{ch:architecture} explains how the
specification tells you to build one.

\subsection{The baseline numbers}

These are the figures every later chapter improves on. Write them down.

\begin{measurebox}{Meridian, before anything is optimised}
\begin{tabular}{@{}lr@{}}
\toprule
\thead{Metric} & \thead{Baseline} \\
\midrule
Wall clock, one task            & 1{,}366.4\,ms \\
Model share of wall clock       & 96.1\,\% \\
Loop iterations                 & 3 \\
Tool calls                      & 3 \\
Total tokens                    & 4{,}009 \\
Cost per task                   & \$0.0125 \\
Tool catalogue                  & 10 tools, 1{,}246 tokens \\
Cache hit rate                  & 0\,\% (nothing cached yet) \\
\bottomrule
\end{tabular}
\end{measurebox}

By Chapter~\ref{ch:optimizing} that task is meaningfully faster and cheaper, and
every step of the change is attributed to a specific technique with a specific
measurement. No hand-waving.

\section{How to read the numbers in this book}
\label{sec:machine}

Two things you are owed up front.

\subsection{The machine}

Every measurement was taken on Apple Silicon (arm64), macOS 15.7.3, Python
3.14.6. Your absolute numbers will differ. The \emph{ratios} are what travel,
and ratios are what the arguments rest on.

Transport figures are loopback figures. That is deliberate. Localhost has
essentially no physics, which isolates protocol overhead so you can see it
without your office wi-fi in the way. Here is what that isolation buys:

\begin{measurebox}{Same call, three transports}
\begin{tabular}{@{}lrr@{}}
\toprule
\thead{Transport} & \thead{Mean} & \thead{Overhead over in-process} \\
\midrule
In-process (control)   & 0.031\,ms & baseline \\
stdio                  & 0.060\,ms & +0.029\,ms \\
Streamable HTTP        & 0.161\,ms & +0.130\,ms \\
\bottomrule
\end{tabular}

\vspace{6pt}
\footnotesize
The in-process transport still serialises and deserialises, so this is a fair
comparison rather than a flattering one.
\end{measurebox}

Subtract the control from the other two and you have transport cost with server
execution removed. That subtraction is the single most useful trick in
Chapter~\ref{ch:measuring}, and it is the only honest way to answer ``is my
server slow, or is my transport slow.''

Add your own network latency on top. A cross-region hop is 40 to 90
milliseconds and swamps all three rows, which is exactly why
Chapter~\ref{ch:transport} cares about connection reuse and
Chapter~\ref{ch:optimizing} cares about round-trip counts.

\subsection{What is simulated}
\label{sec:simulated}

Model inference is simulated. A book cannot ship a reproducible language model,
so Meridian uses a stub with a stated distribution:

\begin{py}
TTFT_MEAN_MS         = 340.0   # time to first token
TTFT_STDEV_MS        = 90.0
MS_PER_OUTPUT_TOKEN  = 7.5
CHARS_PER_TOKEN      = 4.0     # rough, and rough is fine for ratios
INPUT_USD_PER_MTOK   = 3.00
OUTPUT_USD_PER_MTOK  = 15.00
\end{py}

Those constants are calibrated against published figures for mid-size hosted
models in mid-2026. They are stated here so you can disagree with them
precisely rather than vaguely. Edit \texttt{meridian/host/model.py}, rerun
\texttt{make bench}, and every number in the book moves with you.

Everything downstream of inference is measured for real: serialisation,
transport, server execution, cache hit rates, round-trip counts, catalogue token
cost, and process cold starts.

\keyidea{When a measurement box in this book gives you a number, the number came
out of a program you can run. When part of it is modelled, the box says so. If
you catch me violating this, open an issue.}

\section{What you should do on Monday}

Before you read another chapter, go and find out three things about the system
you already have.

\begin{enumerate}[leftmargin=1.6em, itemsep=4pt]
\item \textbf{How many loop iterations does a typical task take?} Not tool
calls. Iterations. Each one is a full model turn. If the answer is more than
four, that is where your latency lives and no amount of server tuning will
touch it.

\item \textbf{How many tokens does your tool catalogue cost per turn?}
Serialise your merged \texttt{tools/list} output, divide the character count by
four, and multiply by the number of turns in a task. Chapter~\ref{ch:tools}
measures a catalogue that costs 4,371 extra tokens on every single turn, which
worked out to \$26 per thousand tasks in pure waste. That was a real catalogue,
and it was not the worst one I have seen.

\item \textbf{What fraction of your requests could have been served from
cache?} If you are not honouring \texttt{ttlMs}, the answer is most of them.
Meridian avoided 577 of 600 resource reads with a 96.2\,\% hit rate on a
realistic access pattern.
\end{enumerate}

Three numbers. An afternoon. They will tell you which chapter to read next
better than my table of contents can.

\section{Where this is going}

The rest of Part~I gets you to the wire: the stateless architecture in
Chapter~\ref{ch:architecture}, both transports byte by byte in
Chapter~\ref{ch:transport}, and OAuth as MCP actually uses it in
Chapter~\ref{ch:authorization}. Part~II is the primitives, one chapter each,
every one of them following the same pattern of mechanics, then cost, then
guidance.

But the frame never changes. Three budgets. Round trips are the enemy. Measure
before you optimise, and measure the thing you actually ship.

Onward.
